\scanpage{10}
\unsure{second|\deleted{second}} element of structure
\[
f:\widetilde{S}\longrightarrow\unsure{\mathscr{J}|\mathcal{J}|\gamma}
\]
These two elements of structure are linked by the condition that $f$ is compatible with the actions of $\pi$, when we let it act on $\mathscr{J}$ via $\alpha:\pi\dashvG(\mathbb{Q})\subset G(\mathbb{R})$ and the \inserted{\uline{right}} action of $G(\mathbb{R})$ on $\mathscr{J}$ (opposite to the left action).

The datum of a pair $(F,\varphi)$ is equivalent to that of a pair $(f,\alpha)$ satisfying the two preceding conditions.

\deleted{One would have to fix \uline{lattices}, and suitable discrete subgroups of $G(\mathbb{Q})$, of type $G(\mathbb{Z})$ commensurable with such a ....}

Another way of presenting the ``things, \unsure{independently|independently} of the choice of a base point $s$ of $S$. \deleted{(\unsure{let us resume|let us start again})}.

\deleted{One says that $F$ is co} Let us note that the fiber functors \inserted{$\Phi_s$ (relative to $\mathbb{Q}$)} on \deleted{Hodge structures} $\mathrm{Hodge}(S)$ defined by the various points $s\in S$ \deleted{\unsure{are}} \inserted{of a \unsure{$\overline{m}$|\deleted{c}} connected component $S^{i}$ of $S$ are} isomorphic ; hence also the $\Phi_s\circ F$ \inserted{$(s\in S^{i})$}. One says that \inserted{the $G$-subordination} $F$ is \uline{compatible with the} \uline{forgetful functor} \uline{on} $C=\operatorname{Rep}(G)$ if the $\Phi_s\circ F$ are isomorphic to the said forgetful functor. Let then, for
% archive sheet number: 9
